\documentclass[../../main.tex]{subfiles}
\begin{document}

\begin{center}
  {\fontsize{48pt}{58pt}\selectfont \textbf{1969}\normalsize\textbf{年度}}
\end{center}

\begin{table}[htb]
  \centering
  \caption{出題テーマと難易度}
  \begin{tabular}{|c|p{9cm}|c|}
    \hline
    \textbf{問題番号} & \textbf{テーマ} & \textbf{難易度} \\
    \hline
    第1問 & 三次式の実一次因子と根の実部 & \\
    \hline
    第2問 & 円内の弦と角・面積の比 & \\
    \hline
    第3問 & 前提が成り立たない命題と対偶 & \\
    \hline
    第4問 & 水槽に流入する溶液の比重 & \\
    \hline
    第5問 & 正六角形の頂点を表すベクトルの期待値 & \\
    \hline
    第6問 & 平方根の和の最小化 & \\
    \hline
  \end{tabular}
\end{table}

\begin{table}[htb]
  \centering
  \caption{解答}
  \renewcommand{\arraystretch}{1.4}
  \begin{tabular}{|c|c|p{9cm}|}
    \hline
    \textbf{問題番号} & \textbf{小問} & \textbf{答え} \\
    \hline
    第1問 & (1) & 実数係数の一次因子をもつ（証明問題） \\
    \cline{2-3}
    & (2) & すべての根の実部が負なら$a,b,c>0$（証明問題） \\
    \hline
    第2問 & -- & $\tan\alpha\tan\beta=\dfrac{\triangle ADB}{\triangle ACB}=\dfrac{DP}{CP}$（証明問題） \\
    \hline
    第3問 & -- &  \\
    \hline
    第4問 & -- & $1+\dfrac{a}{bk}(1-e^{-bk})$ \\
    \hline
    第5問 & -- & $a^2+1$ \\
    \hline
    第6問 & -- & $\min\{a,1\}$ \\
    \hline
  \end{tabular}
\end{table}

\paragraph{試験形態}
6問．試験時間は要確認．

\end{document}
